\documentclass[aspectratio=169,11pt]{beamer}
\usepackage[utf8]{inputenc}
\usepackage[T1]{fontenc}
\usepackage[spanish,es-noquoting]{babel}
\usepackage{tikz,booktabs}
\usepackage[absolute,overlay]{textpos}
\usetikzlibrary{shapes.geometric,arrows.meta,positioning,calc}

\usetheme{default}
\definecolor{cmAzul}{RGB}{20,58,120}
\definecolor{cmVerde}{RGB}{28,150,110}
\definecolor{cmAmbar}{RGB}{200,120,30}
\definecolor{cmGris}{RGB}{95,105,120}
\definecolor{cmRojo}{RGB}{170,45,45}

\colorlet{itmazul}{cmAzul}
\colorlet{acento}{cmVerde}
\colorlet{verde}{cmVerde}
\colorlet{rojo}{cmAmbar}

% ============================================================
% Logo ComunicaMedellin
% ============================================================
% Concepto: una senal de difusion saliendo de un punto.
% Los tres arcos crecientes son las tres zonas del fan-out;
% el punto solido es la alerta que se emite una sola vez.
% ============================================================

\newcommand{\logomarca}[1][1]{%
\begin{tikzpicture}[scale=#1]
  % ondas de difusion
  \draw[line width=1.1pt,cmVerde,line cap=round]
    (0.16,0) arc (0:62:0.16) (0.16,0) arc (0:-62:0.16);
  \draw[line width=1.1pt,cmAzul!75,line cap=round]
    (0.30,0) arc (0:58:0.30) (0.30,0) arc (0:-58:0.30);
  \draw[line width=1.1pt,cmAzul!45,line cap=round]
    (0.44,0) arc (0:54:0.44) (0.44,0) arc (0:-54:0.44);
  % punto emisor
  \fill[cmAzul] (0,0) circle (0.075);
\end{tikzpicture}}

% Logo con el nombre al lado
\newcommand{\logotipo}[1][1]{%
\begin{tikzpicture}[scale=#1]
  \node[inner sep=0pt] at (0,0) {\logomarca[#1]};
  \node[right=2pt of {(0.5,0)},anchor=west,inner sep=0pt] at (0.52,0)
    {\fontsize{#1*11}{#1*12}\selectfont
     \textcolor{cmAzul}{\textbf{Comunica}}\textcolor{cmVerde}{\textbf{Medellín}}};
\end{tikzpicture}}

\setbeamercolor{structure}{fg=itmazul}
\setbeamercolor{frametitle}{fg=white,bg=itmazul}
\setbeamertemplate{navigation symbols}{}
\setbeamercolor{title}{fg=itmazul}
\setbeamercolor{itemize item}{fg=acento}
\setbeamercolor{itemize subitem}{fg=acento}
\setbeamercolor{block title}{fg=white,bg=itmazul}
\setbeamercolor{block body}{bg=itmazul!6}
\setbeamertemplate{footline}{%
  \hspace{6pt}\raisebox{1pt}{\logomarca[0.5]}%
  \hfill\tiny\textcolor{cmGris}{\insertframenumber}\hspace{8pt}\vspace{3pt}}
\setbeamerfont{frametitle}{size=\large}

% Referencia al material del curso, con modulo explicito
\newcommand{\mI}[1]{\textcolor{acento}{\small\textbf{[M1,~Sec.~#1]}}}

\newcommand{\quien}[1]{%
  \begin{textblock*}{4cm}(8.6cm,0.32cm)
    \tiny\textcolor{white}{\textbf{#1}}
  \end{textblock*}}

% Etiqueta de dato medido en el cluster
\newcommand{\medido}[1]{\textcolor{cmVerde!70!black}{\textbf{#1}}}

\title{\textbf{Modelado de Datos en Cassandra}}
\subtitle{ComunicaMedellín — Fan-out, hotspots y consistencia}
\author{Grupo 3}
\institute{Tópicos Avanzados de Bases de Datos · Módulo 2 · ITM}
\date{2026}

\begin{document}

% ═════════════════════════════════════════════ PORTADA
\begin{frame}[plain]
\vfill
\centering
{\logomarca[3.2]}\\[14pt]
{\Huge\textcolor{cmAzul}{\textbf{Comunica}}\textcolor{cmVerde}{\textbf{Medellín}}}\\[8pt]
{\large\textcolor{cmGris}{Del «qué motor» al «qué modelo».}}\\[20pt]
\rule{0.35\textwidth}{0.4pt}\\[12pt]
{\footnotesize\textcolor{cmGris}{Tópicos Avanzados de Bases de Datos · Módulo 2 · ITM}}\\[6pt]
{\footnotesize\textbf{Grupo 3} \quad
 \textcolor{cmGris}{Julián · Elizabeth · Simón · Nilton · Lina}}
\vfill
\end{frame}

% ═════════════════════════════════════════════ 1. DE DÓNDE VENIMOS
\begin{frame}{Dónde quedamos y qué falta}
\quien{Julián}
\vspace{2pt}
\centering
\begin{tikzpicture}[font=\scriptsize,
  hecho/.style={draw=cmVerde!50,fill=cmVerde!8,rounded corners=3pt,
                minimum height=13mm,text width=38mm,align=center,inner sep=4pt},
  falta/.style={draw=cmAmbar!60,fill=cmAmbar!8,rounded corners=3pt,
                minimum height=13mm,text width=38mm,align=center,inner sep=4pt},
  fl/.style={-{Stealth[length=2mm]},cmAzul!55,thick}]
\node[hecho] (m1) at (0,0)
  {\textbf{Módulo 1} — \emph{qué motor}\\[2pt]
   \tiny Cassandra AP · Kafka · ClickHouse\\
   \tiny 5 cargas de trabajo identificadas};
\node[falta] (m2) at (5.6,0)
  {\textbf{Módulo 2} — \emph{qué modelo}\\[2pt]
   \tiny tablas, claves, consistencia\\
   \tiny y qué pasa cuando algo falla};
\draw[fl] (m1) -- (m2);
\end{tikzpicture}

\vspace{10pt}
\footnotesize
\begin{minipage}{0.88\textwidth}
Elegir Cassandra no resuelve nada por sí solo. Falta lo difícil:
\textbf{cómo se guardan realmente los datos} para que el sistema responda rápido,
no se sature al enviar una alerta a 100.000 personas, y se recupere cuando un
nodo se cae.

\vspace{7pt}
\begin{center}
\begin{tabular}{@{}p{5.0cm}p{5.6cm}@{}}
\toprule
\textbf{Lo que trae el Módulo 2} & \textbf{Cómo lo demostramos} \\
\midrule
7 tablas con clave justificada & consultas reales, medidas \\
Fan-out de 100.000 registros & 10.000 entregas en vivo \\
Duplicados de Kafka & reenviamos el mismo evento \\
Nodo caído mucho tiempo & lo tumbamos y lo reparamos \\
\bottomrule
\end{tabular}
\end{center}
\end{minipage}
\end{frame}

% ═════════════════════════════════════════════ 2. QUERY-FIRST
\begin{frame}{El modelo sale de las consultas}
\quien{Julián}
\footnotesize
\vspace{2pt}
\begin{columns}[T]
\begin{column}{0.47\textwidth}
\begin{center}\textbf{\textcolor{cmGris}{Relacional}}\end{center}
\vspace{-2pt}
\begin{tikzpicture}[font=\tiny,
  c/.style={draw=cmGris!45,fill=cmGris!6,rounded corners=2pt,
            minimum width=30mm,minimum height=6.5mm,align=center},
  f/.style={-{Stealth[length=1.5mm]},cmGris!60}]
\node[c] (a) at (0,0)    {modelo las \textbf{entidades}};
\node[c] (b) at (0,-1.0) {normalizo (3NF)};
\node[c] (c) at (0,-2.0) {consulto \textbf{lo que quiera}};
\draw[f] (a)--(b); \draw[f] (b)--(c);
\end{tikzpicture}
\end{column}
\begin{column}{0.47\textwidth}
\begin{center}\textbf{\textcolor{cmVerde}{Cassandra}}\end{center}
\vspace{-2pt}
\begin{tikzpicture}[font=\tiny,
  c/.style={draw=cmVerde!50,fill=cmVerde!8,rounded corners=2pt,
            minimum width=30mm,minimum height=6.5mm,align=center},
  f/.style={-{Stealth[length=1.5mm]},cmVerde!65}]
\node[c] (a) at (0,0)    {listo las \textbf{consultas}};
\node[c] (b) at (0,-1.0) {una tabla por consulta};
\node[c] (c) at (0,-2.0) {\textbf{duplico} el dato si hace falta};
\draw[f] (a)--(b); \draw[f] (b)--(c);
\end{tikzpicture}
\end{column}
\end{columns}

\vspace{9pt}
\begin{minipage}{0.93\textwidth}\footnotesize
\textbf{Por qué al revés:} los datos se reparten por el \emph{hash} de la clave de
partición \mI{3.3}. Solo se puede buscar rápido \textbf{por esa clave} —así que
primero hay que saber qué se va a preguntar.

\vspace{5pt}
\textcolor{cmAmbar}{\textbf{Al escribir nuestras 5 consultas, dos no tenían
tabla que las respondiera.}} La solución en Cassandra no es un índice: es
\textbf{otra tabla}. Por eso pasamos de 3 tablas a 7.
\end{minipage}
\end{frame}

% ═════════════════════════════════════════════ 3. LAS 5 CONSULTAS
\begin{frame}{Las cinco consultas reales}
\quien{Julián}
\vspace{2pt}
\scriptsize
\begin{center}
\begin{tabular}{@{}p{0.6cm}p{5.0cm}p{2.0cm}p{4.2cm}@{}}
\toprule
& \textbf{Qué necesita el sistema} & \textbf{Frecuencia} & \textbf{Tabla que la sirve} \\
\midrule
Q1 & Presencia actual de un usuario & $\sim$10.000/s & \texttt{presencia} \\
Q2 & Historial de alertas de un usuario & $\sim$500/s & \texttt{historial\_usuario} \\
Q3 & Alertas en un rango de fechas & $\sim$100/s & \texttt{historial\_usuario} \\
Q4 & \textcolor{cmAmbar}{Usuarios en línea de una zona} & $\sim$20/día & \textcolor{cmAmbar}{\texttt{presencia\_por\_comuna}} \\
Q5 & \textcolor{cmAmbar}{Quién tiene pendiente una alerta} & $\sim$50/día & \textcolor{cmAmbar}{\texttt{estado\_entrega}} \\
\bottomrule
\end{tabular}
\end{center}

\vspace{8pt}
\footnotesize
\begin{columns}[T]
\begin{column}{0.48\textwidth}
\textbf{\textcolor{cmAmbar}{Las dos que faltaban}}\\[3pt]
\scriptsize
\textbf{Q4} no se puede responder desde \texttt{presencia}: está particionada por
\texttt{usuario\_id} y la zona no participa en el \emph{hash}.\\[3pt]
\textbf{Q5} sencillamente no se guardaba en ningún lado.
\end{column}
\begin{column}{0.48\textwidth}
\textbf{\textcolor{cmVerde}{Resultado medido}}\\[3pt]
\scriptsize
Las cinco resuelven en \medido{menos de 4 ms} (p95) porque cada una toca
\textbf{una sola partición}.\\[3pt]
Ninguna usa \texttt{ALLOW FILTERING}, ninguna ordena en memoria, ninguna cruza
tablas.
\end{column}
\end{columns}
\end{frame}

% ═════════════════════════════════════════════ 4. EL MODELO
\begin{frame}{El modelo: una tabla por pregunta}
\quien{Elizabeth}
\vspace{1pt}
\centering
\begin{tikzpicture}[font=\tiny,scale=1.12,every node/.style={transform shape},
  t/.style={draw=cmAzul!45,fill=white,rounded corners=2pt,
            text width=40mm,align=left,inner sep=3pt},
  q/.style={draw=cmVerde!55,fill=cmVerde!8,rounded corners=5pt,
            text width=15mm,align=center,inner sep=2.5pt},
  k/.style={draw=cmVerde!55,fill=cmVerde!10,rounded corners=2pt,
            text width=34mm,align=center,inner sep=3pt},
  l/.style={-{Stealth[length=1.3mm]},cmVerde!60},
  e/.style={-{Stealth[length=1.5mm]},cmAzul!60,thick}]

\node[k] (kaf) at (-1.0,1.15) {\textbf{Kafka} — 1 alerta\\\textcolor{cmGris}{partición por zona}};

\node[q] (q1) at (-3.85,0.15)  {\textbf{Q1} presencia};
\node[q] (q2) at (-3.85,-0.85) {\textbf{Q2--Q3} historial};
\node[q] (q3) at (-3.85,-1.85) {\textbf{Q4} por zona};
\node[q] (q4) at (-3.85,-2.85) {\textbf{Q5} entrega};

\node[t] (t1) at (-0.1,0.15)  {\texttt{presencia}\\\textcolor{cmAzul}{PK} \texttt{(usuario\_id)} · \textcolor{cmAmbar}{\texttt{ONE}}};
\node[t] (t2) at (-0.1,-0.85) {\texttt{historial\_usuario}\\\textcolor{cmAzul}{PK} \texttt{((usuario\_id), enviado\_en}$\downarrow$\texttt{)}};
\node[t] (t3) at (-0.1,-1.85) {\texttt{presencia\_por\_comuna}\\\textcolor{cmAzul}{PK} \texttt{((comuna, estado), usuario)}};
\node[t] (t4) at (-0.1,-2.85) {\texttt{estado\_entrega}\\\textcolor{cmAzul}{PK} \texttt{((alerta\_id, bucket), usuario)}};

\node[t] (t5) at (4.15,1.15)  {\texttt{alertas}\\\textcolor{cmAzul}{PK} \texttt{(alerta\_id)} · catálogo};
\node[t] (t6) at (4.15,-1.35) {\texttt{eventos\_procesados}\\\textcolor{cmAzul}{PK} \texttt{(evento\_id)} · guardia \texttt{LWT}};

\foreach \i in {1,2,3,4} { \draw[l] (q\i.east) -- (t\i.west); }
\draw[e] (kaf.east) -- (t5.west) node[midway,above,font=\tiny,text=cmGris] {1 escr.};
\draw[e] (kaf.south) -- ++(0,-0.35) -| (t1.north);
\draw[e] (kaf.west) -- ++(-1.75,0) |- (t2.west);
\node[font=\tiny,text=cmGris,align=center,fill=white,inner sep=1pt]
  at (-2.85,0.62) {fan-out\\100.000};
\draw[e,dashed] (t2.east) -- (t6.west);
\end{tikzpicture}

\vspace{4pt}
\scriptsize
\textcolor{cmGris}{Verde = ruta de lectura · Azul = ruta de escritura.}
\textbf{\texttt{presencia\_por\_comuna} es la misma información que
\texttt{presencia}, indexada al revés} \mI{2.2}.
\end{frame}

% ═════════════════════════════════════════════ 5. CONSISTENCIA
\begin{frame}{La consistencia se decide por consulta}
\quien{Elizabeth}
\footnotesize
\vspace{2pt}
Lo que se configura \textbf{una sola vez} es la replicación:
\texttt{RF=3} como \texttt{\{medellin:2, miami:1\}}.
Lo que se elige \textbf{en cada operación} es el nivel de consistencia \mI{1.2}.

\vspace{7pt}
\begin{columns}[T]
\begin{column}{0.44\textwidth}
\centering\scriptsize
\begin{tabular}{@{}llc@{}}
\toprule
\textbf{Escritura} & \textbf{Lectura} & \textbf{Garantía} \\
\midrule
\texttt{ONE} & \texttt{ONE} & eventual \\
\texttt{QUORUM} & \texttt{ONE} & durable \\
\texttt{QUORUM} & \texttt{QUORUM} & \textcolor{cmVerde}{\textbf{inmediata}} \\
\bottomrule
\end{tabular}\\[6pt]
\textcolor{cmAzul}{\textbf{$R + W > RF$}}\\
\tiny\textcolor{cmGris}{con RF=3, el quórum son 2 réplicas}
\end{column}
\begin{column}{0.53\textwidth}
\scriptsize
\textbf{En nuestro sistema, a propósito:}\\[4pt]
\begin{tabular}{@{}p{2.3cm}p{3.4cm}@{}}
\toprule
\texttt{presencia} & \texttt{ONE} — que llegue 2 s tarde no cambia ninguna decisión \\
\addlinespace
\texttt{historial} & \texttt{QUORUM} escr. — es seguridad pública, no puede perderse \\
\bottomrule
\end{tabular}\\[6pt]
\textcolor{cmAmbar}{\textbf{Ojo:}} el historial escribe con \texttt{QUORUM} pero
lee con \texttt{ONE}. Necesitamos \emph{durabilidad}, no \emph{lectura
inmediata}: que el usuario la vea 300 ms después es irrelevante.
\end{column}
\end{columns}

\vspace{7pt}
\centering\scriptsize
\begin{tabular}{@{}lccc@{}}
\toprule
\textbf{Medido sobre la misma consulta} & \texttt{ONE} & \texttt{LOCAL\_QUORUM} & \texttt{QUORUM} \\
\midrule
Latencia p50 & \medido{1,13 ms} & 1,38 ms & 1,54 ms \\
Sobrecosto & — & $+21{,}7$~\% & $+36{,}6$~\% \\
\bottomrule
\end{tabular}\\[3pt]
\tiny\textcolor{cmGris}{En red local la diferencia es de 0,41 ms. Con Miami real
(60--80 ms) \texttt{LOCAL\_QUORUM} deja de ser un detalle.}
\end{frame}

% ═════════════════════════════════════════════ 6. FAN-OUT
\begin{frame}{Fan-out: 1 alerta, 100.000 escrituras}
\quien{Simón}
\footnotesize
\vspace{2pt}
\begin{columns}[T]
\begin{column}{0.46\textwidth}
\textbf{Lo que cuesta en disco}\\[4pt]
\scriptsize
\begin{tabular}{@{}lrr@{}}
\toprule
& \textbf{Filas} & \textbf{RF=3} \\
\midrule
1 alerta & 100 K & 6,44 MB \\
1 día (20 alertas) & 2 M & 0,13 GB \\
1 mes & 60 M & 3,77 GB \\
1 año & 730 M & \medido{45,9 GB} \\
\bottomrule
\end{tabular}\\[5pt]
\textcolor{cmVerde}{\scriptsize\textbf{15 GB por nodo: irrelevante.}}\\
\tiny\textcolor{cmGris}{El fan-out no es un problema de espacio,
sino de \textbf{tasa de escritura y distribución}.}
\end{column}
\begin{column}{0.51\textwidth}
\textbf{Lo que cuesta en tiempo} \tiny\textcolor{cmGris}{(medido con 10.000)}\\[4pt]
\scriptsize
\begin{tabular}{@{}lr@{\hspace{6mm}}lr@{}}
\toprule
\multicolumn{2}{@{}l}{\textbf{Caudal}} & \multicolumn{2}{l}{\textbf{Latencia}} \\
\midrule
Entregas & 9.999 & media & 90,8 ms \\
Tiempo & 1,33 s & p50 & 93,4 ms \\
Global & \medido{7.509/s} & p95 & 142,5 ms \\
vs serie & 2,5$\times$ & p99 & \medido{156,4 ms} \\
\bottomrule
\end{tabular}\\[5pt]
\tiny\textcolor{cmGris}{\textbf{No son lo mismo:} el caudal mide el sistema, la
latencia mide lo que espera \emph{una} escritura. Con 400 en vuelo el caudal sube
y la latencia individual empeora.}
\end{column}
\end{columns}

\vspace{7pt}
\begin{minipage}{0.95\textwidth}\scriptsize
\textcolor{cmAmbar}{\textbf{Consecuencia de diseño:}} a 7.509/s el fan-out completo
tarda $\sim$13 s. Por eso \textbf{es asíncrono}: la Alcaldía recibe confirmación
cuando la alerta entra en Kafka \mI{4.2}, no cuando terminan las 100.000
escrituras. Si llegan 10 alertas a la vez, Kafka amortigua: se retrasa la entrega,
no se pierde nada.
\end{minipage}
\end{frame}

% ═════════════════════════════════════════════ 7. HOTSPOT
\begin{frame}{El error que parece correcto}
\quien{Simón}
\footnotesize
\vspace{2pt}
\centering
\begin{tikzpicture}[font=\tiny,
  mal/.style={draw=cmRojo!55,fill=cmRojo!7,rounded corners=2pt,
              minimum width=6mm,minimum height=6mm},
  bien/.style={draw=cmVerde!55,fill=cmVerde!10,rounded corners=1.5pt,
               minimum width=3.4mm,minimum height=3.4mm}]
% izquierda: 1 particion gigante
\node[font=\scriptsize,text=cmRojo] at (-3.2,1.25) {\textbf{PK (alerta\_id)}};
\node[mal,minimum width=17mm,minimum height=17mm] at (-3.2,0) {};
\node[font=\scriptsize,text=cmRojo] at (-3.2,0) {\textbf{1}};
\node[text=cmGris] at (-3.2,-1.35) {1 partición · 1.909 KB};
\node[text=cmRojo!85] at (-3.2,-1.72) {toda la carga a 1 punto del anillo};

\node[font=\Large,text=cmGris] at (-0.6,0) {vs};

% derecha: 32 particiones
\node[font=\scriptsize,text=cmVerde!65!black] at (2.6,1.25) {\textbf{PK (alerta\_id, bucket)}};
\foreach \i in {0,...,31}{
  \node[bien] at ({1.55+mod(\i,8)*0.42},{0.62-floor(\i/8)*0.42}) {};
}
\node[text=cmGris] at (2.6,-1.35) {32 particiones · 71,7 KB c/u};
\node[text=cmVerde!65!black] at (2.6,-1.72) {repartidas por todo el anillo};
\end{tikzpicture}

\vspace{8pt}
\begin{minipage}{0.93\textwidth}\scriptsize
\textcolor{cmAmbar}{\textbf{Lo contraintuitivo:}} 3,73 MB \emph{caben de sobra} en
el límite de 100 MB por partición. \textbf{El problema no es el tamaño.} Es que las
100.000 escrituras van a \emph{un} \emph{token} del anillo, y por tanto siempre a
las mismas 3 réplicas.

\vspace{4pt}
Con 3 nodos no se nota, porque todos son réplica de todo. \textbf{Con 12 nodos:
3 saturados y 9 ociosos.} Además, reparar y compactar una partición gigante es una
operación atómica costosa; 32 pequeñas se paralelizan.

\vspace{4pt}
\textcolor{cmVerde!65!black}{\textbf{La clave de partición no se elige por la
entidad, sino por cómo distribuye la carga.}}
\end{minipage}
\end{frame}

% ═════════════════════════════════════════════ 8. IDEMPOTENCIA
\begin{frame}{¿Llega la alerta dos veces?}
\quien{Nilton}
\footnotesize
\vspace{2pt}
Kafka entrega \textbf{al menos una vez}: si un \emph{worker} muere antes de
confirmar el \emph{offset}, al reiniciar reprocesa el mensaje. En un fan-out eso
es reenviar la alerta a 100.000 personas.

\vspace{7pt}
\centering
\begin{tabular}{@{}p{4.6cm}p{4.6cm}@{}}
\toprule
\textcolor{cmVerde!65!black}{\textbf{Timestamp del EVENTO}} &
\textcolor{cmRojo}{\textbf{Timestamp del WORKER}} \\
\midrule
\scriptsize El productor lo fija una vez y viaja dentro del mensaje &
\scriptsize Cada reintento llama a \texttt{now()} \\
\addlinespace
\scriptsize Misma clave primaria $\Rightarrow$ \emph{last-write-wins} &
\scriptsize Clave distinta $\Rightarrow$ fila nueva \\
\addlinespace
\hfill\medido{1 fila} tras 3 reenvíos\hfill\null &
\hfill\textcolor{cmRojo}{\textbf{3 filas}} tras 3 reenvíos\hfill\null \\
\bottomrule
\end{tabular}

\vspace{8pt}
\begin{minipage}{0.93\textwidth}\scriptsize
\textcolor{cmVerde!65!black}{\textbf{La defensa no cuesta nada.}} Cassandra no
tiene \texttt{UNIQUE}: un \texttt{INSERT} sobre una clave existente es un
\emph{upsert} silencioso. Eso, que parece una carencia, \textbf{es} la garantía —
basta con que la clave primaria derive del evento y no del momento de procesarlo.
Basta 1 ms de diferencia para romperlo.

\vspace{4pt}
\textcolor{cmAmbar}{\textbf{Lo que sí hay que proteger:}} el \emph{push} al celular
ya salió y no se retira. Para eso el \emph{worker} reclama el evento con
\texttt{IF NOT EXISTS} (Paxos, 10--20$\times$ más lento) \textbf{solo antes del
push} — nunca en las 100.000 escrituras.
\end{minipage}
\end{frame}

% ═════════════════════════════════════════════ 9. NODO CAÍDO
\begin{frame}{Un nodo se cae y tarda en volver}
\quien{Nilton}
\footnotesize
\vspace{2pt}
\begin{columns}[T]
\begin{column}{0.45\textwidth}
\textbf{Tres mecanismos, en orden}\\[4pt]
\scriptsize
\begin{tabular}{@{}p{2.35cm}p{2.5cm}@{}}
\toprule
\textbf{Hinted handoff} & automático, solo \textbf{3 h} \\
\addlinespace
\textbf{Read repair} & solo lo que alguien lee \\
\addlinespace
\textbf{\texttt{nodetool}} \textbf{\texttt{repair}} & \textcolor{cmVerde!65!black}{única garantía completa} \\
\bottomrule
\end{tabular}\\[6pt]
\textcolor{cmAmbar}{\textbf{Dos umbrales:}}\\
\tiny \textbf{3 h} — pasa la ventana de \emph{hints}: hay que reparar.\\
\tiny \textbf{10 días} (\texttt{gc\_grace}) — los borrados \textbf{resucitan}:
hay que reemplazar el nodo, no reintegrarlo.
\end{column}
\begin{column}{0.52\textwidth}
\textbf{Medido en vivo} \tiny\textcolor{cmGris}{(2.000 particiones)}\\[4pt]
\scriptsize
\begin{tabular}{@{}lccc@{}}
\toprule
\textbf{Particiones} & \texttt{cass1} & \texttt{cass2} & \texttt{cass3} \\
\midrule
línea base & 0 & 0 & 0 \\
tras escribir sin \texttt{cass2} & 1.999 & \textcolor{cmRojo}{\textbf{0}} & 1.999 \\
tras \texttt{repair} & 1.999 & \medido{1.999} & 1.999 \\
\bottomrule
\end{tabular}\\[6pt]
\tiny\textcolor{cmGris}{El clúster \textbf{siguió aceptando escrituras} con un
nodo caído: eso es la A de AP \mI{1.2}. Con CP las habría rechazado.}
\end{column}
\end{columns}

\vspace{7pt}
\begin{minipage}{0.95\textwidth}\scriptsize
\textcolor{cmAmbar}{\textbf{Lo que descubrimos midiendo:}} el hueco
\textbf{no se ve con un \texttt{SELECT}}. Con RF=3 sobre 3 nodos los tres son
réplica de todo, así que una lectura con \texttt{ONE} dirigida a \texttt{cass2} la
sirve cualquier réplica viva y devuelve el dato correcto.

\vspace{3pt}
\textcolor{cmVerde!65!black}{\textbf{Consistencia eventual no significa que el
cliente lea mal: significa que las réplicas divergen por dentro mientras el
servicio responde bien.}} Por eso hay que reparar aunque nada parezca fallar.
\end{minipage}
\end{frame}

% ═════════════════════════════════════════════ 10. OLTP vs OLAP
\begin{frame}{La consulta que no le toca a Cassandra}
\quien{Lina}
\footnotesize
\vspace{2pt}
\centering
\textbf{«¿Cuántos usuarios únicos hay en línea, por zona?»}\\
\tiny\textcolor{cmGris}{Una agregación sobre una columna que no es clave de partición.}

\vspace{7pt}
\begin{columns}[T]
\begin{column}{0.47\textwidth}
\scriptsize\textbf{\textcolor{cmRojo}{En Cassandra}}\\[3pt]
\tiny
\texttt{SELECT zona, count(*)}\\
\texttt{FROM presencia GROUP BY zona;}\\[4pt]
\textcolor{cmRojo}{\textbf{$\rightarrow$ RECHAZADA}}\\
\tiny\textcolor{cmGris}{solo agrupa por columnas de la clave primaria;
las filas de una zona están por todo el anillo}\\[4pt]
Forzada con \texttt{ALLOW FILTERING} responde, pero el \texttt{GROUP BY} lo
termina haciendo \textbf{el cliente en memoria}.
\end{column}
\begin{column}{0.47\textwidth}
\scriptsize\textbf{\textcolor{cmVerde!65!black}{En ClickHouse}}\\[3pt]
\tiny
\texttt{SELECT zona, uniqExact(usuario\_id)}\\
\texttt{FROM presencia\_eventos}\\
\texttt{GROUP BY zona;}\\[4pt]
\textcolor{cmVerde!65!black}{\textbf{$\rightarrow$ es su caso de uso}}\\
\tiny\textcolor{cmGris}{columnar: lee 2 de 6 columnas, comprimidas con
\emph{dictionary encoding} porque zona tiene 3 valores}
\end{column}
\end{columns}

\vspace{7pt}
\centering\scriptsize
\begin{tabular}{@{}lrrr@{}}
\toprule
\textbf{Motor} & \textbf{Filas} & \textbf{Latencia} & \textbf{Costo por fila} \\
\midrule
Cassandra (\texttt{ALLOW FILTERING}) & 10.000 & 77,9 ms & 7,79 µs \\
ClickHouse (columnar) & 2.000.000 & 77,9 ms & \medido{0,04 µs} \\
\bottomrule
\end{tabular}

\vspace{5pt}
\begin{minipage}{0.93\textwidth}\scriptsize
\textcolor{cmAmbar}{\textbf{Cuidado al leer el dato:}} la latencia es la misma,
pero \textbf{no sobre el mismo volumen} — ClickHouse agregó 200 veces más filas en
ese tiempo. Comparar latencias directas \emph{subestima} la diferencia.

\vspace{3pt}
No es que un motor sea mejor: \textbf{resuelven preguntas distintas}, y por eso la
arquitectura tiene los dos \mI{2.2}.
\end{minipage}
\end{frame}

% ═════════════════════════════════════════════ 11. CONCLUSIONES
\begin{frame}{Lo que nos llevamos}
\quien{Lina}
\footnotesize
\vspace{3pt}
\begin{enumerate}[<+->]\itemsep5pt
\item \textbf{Las consultas produjeron el modelo, no al revés.} Escribirlas antes
de diseñar reveló dos que ninguna tabla podía servir. Duplicar datos es el precio
explícito de que las cinco respondan en \medido{menos de 4 ms}.

\item \textbf{La clave de partición se elige por distribución, no por entidad.}
Particionar por \texttt{alerta\_id} es lo natural desde el modelo de entidades —y
es la decisión equivocada: \medido{1.909 KB en 1 partición} frente a
\medido{71,7 KB en 32}.

\item \textbf{La consistencia es propiedad de cada operación.} La replicación se
configura una vez; el nivel se decide por consulta, bajo $R+W>RF$.

\item \textbf{La idempotencia se diseña, no se añade.} Que la clave primaria
derive del evento hace que reprocesar Kafka sea gratis.

\item \textbf{Hay preguntas que el modelo no debe responder.} Reconocer el límite
del motor operacional es parte del diseño, no una carencia.
\end{enumerate}
\end{frame}

% ═════════════════════════════════════════════ CIERRE
\begin{frame}[plain]\centering\vfill
{\logomarca[2.2]}\\[12pt]
{\Large\textcolor{cmAzul}{\textbf{¿Preguntas?}}}\\[10pt]
{\footnotesize\textcolor{cmGris}{Grupo 3 · Módulo 2 · ITM}}\\[14pt]
{\tiny\textcolor{cmGris}{\texttt{make demo-fanout} · \texttt{make demo-idempotencia} ·
\texttt{make demo-reparacion} · \texttt{make demo-analitica}}}
\vfill\end{frame}

\end{document}
